% ======================================================================
% CAPÍTULO 7 — CONCLUSÃO
% ======================================================================
\chapter{CONCLUSÃO}

Este capítulo retoma o problema e os objetivos propostos na introdução, sintetiza
as contribuições do trabalho, apresenta as limitações e indica direções para
trabalhos futuros.

\section{Síntese do trabalho}

Este trabalho apresentou o projeto e a implementação de um sistema de automação
e supervisão para a preparação de amostras em especiação de mercúrio, concebido
para substituir um sistema legado baseado em LabVIEW. A solução adota uma
arquitetura mestre-escravo composta por um Raspberry Pi (lógica de alto nível em
Python), um Arduino Uno (aquisição de dados em tempo real) e uma IHM Web em
React e TypeScript, integrados por um protocolo JSON a 4~Hz.

A hipótese de trabalho foi confirmada no nível de validação alcançado: a
arquitetura proposta atende aos requisitos funcionais e de segurança definidos,
proporcionando controle sequenciado do processo por máquina de estados finita,
controle PID misto da temperatura, persistência auditável dos parâmetros e
parada segura automática por watchdog duplo. A validação por testes
automatizados --- 67 testes aprovados distribuídos entre backend, firmware e IHM
Web --- demonstrou a correção da lógica de controle, a integração entre as
camadas e a efetividade dos mecanismos de segurança no nível de software.

\section{Atendimento aos objetivos}

Retomando os objetivos específicos estabelecidos no Capítulo~1:

\begin{enumerate}[label=\arabic*)]
  \item O processo foi modelado em uma FSM com os estados \texttt{SAFE},
  \texttt{T0\_DERIV} a \texttt{T3\_PURGA} e \texttt{MANUAL}, e a matriz de
  acionamento dos atuadores foi implementada e validada por testes;
  \item O controle PID misto do Tubo U --- razão de taxas abaixo de
  $0\,\celsius$ e PID em malha fechada acima --- e o controle do Forno 2 a
  $700\,\celsius$ foram implementados e verificados por testes de unidade;
  \item O protocolo JSON estruturado entre Raspberry Pi e Arduino, com handshake
  de parametrização e pacotes de escrita e leitura a 4~Hz, foi definido e
  implementado em ambas as plataformas;
  \item O firmware do DAQ foi desenvolvido com leitura de termopares via SPI,
  acionamento de atuadores (incluindo o tratamento da bomba por pulso) e
  watchdog de segurança;
  \item A IHM Web foi desenvolvida com monitoramento em tempo real, Diagrama de
  Tempos, gráficos de tendência e painéis de controle e de configuração;
  \item A persistência atômica dos parâmetros do método em JSON, com backup
  rotativo e validação de faixas, foi implementada e testada;
  \item A validação foi conduzida por testes unitários, de integração,
  funcionais e de segurança, seguindo a estratégia de desenvolvimento orientado
  a testes.
\end{enumerate}

Dessa forma, o objetivo geral --- projetar e implementar um sistema de automação
e supervisão para a preparação de amostras em especiação de mercúrio, com
parâmetros persistentes e segurança operacional --- foi alcançado no âmbito do
escopo definido.

\section{Contribuições}

As principais contribuições deste trabalho são:

\begin{itemize}
  \item a modelagem explícita e verificável do sequenciamento do processo em uma
  máquina de estados finita, com matriz de acionamento documentada;
  \item a estratégia de controle misto da rampa do Tubo U, que contorna a
  limitação de leitura do termopar na região criogênica pela razão de taxas e
  emprega PID em malha fechada na região de leitura confiável;
  \item a arquitetura mestre-escravo de baixo custo com separação entre a lógica
  de supervisão e a execução em tempo real;
  \item o protocolo de comunicação JSON aberto e legível, facilitando
  depuração, manutenção e evolução;
  \item o projeto de segurança em camadas, com estado seguro por padrão e
  watchdog duplo, suprindo uma lacuna do sistema legado;
  \item a validação sistemática por testes automatizados, conferindo
  rastreabilidade entre requisitos, implementação e resultados.
\end{itemize}

\section{Limitações}

As limitações deste trabalho decorrem, principalmente, da indisponibilidade da
bancada física durante o desenvolvimento. A validação foi realizada em ambiente
simulado; os critérios de desempenho térmico --- linearidade da rampa e
estabilidade do Forno 2 --- e a calibração da curva de aquecimento do Tubo U
permanecem pendentes de verificação experimental. Registra-se, ainda, que
questões como a identificação exata do conversor de termopar e a definição da
margem de proteção de temperatura dependem da documentação e das decisões do
responsável técnico do equipamento.

\section{Trabalhos futuros}

Como desdobramentos naturais deste trabalho, sugere-se:

\begin{enumerate}[label=\arabic*)]
  \item validar o sistema em bancada com o processo físico, calibrando a curva
  de aquecimento do Tubo U e verificando a linearidade da rampa e a estabilidade
  do Forno 2;
  \item implementar a integração automática com o detector Lumex e a exportação
  das curvas geradas (CSV/PDF) para registro e auditoria;
  \item adicionar autenticação de operadores e trilha de auditoria das ações;
  \item desenvolver um simulador de processo para treinamento de operadores;
  \item implementar a sintonização automática (auto-tune) dos controladores PID;
  \item avaliar a substitui\c{c}\~ao do bit-bang SPI por hardware SPI nativo e a
  inclus\~ao de medi\c{c}\~ao redundante de temperatura para maior robustez.
\end{enumerate}

Em síntese, o trabalho atingiu o objetivo de modernizar a automação e a
supervisão do processo de preparação de amostras para especiação de mercúrio,
entregando uma solução modular, segura e validada por testes, que estabelece uma
base sólida para a validação experimental e para a evolução contínua do sistema.
